Der Organisator eines Weiterbildungsevents einer grossen Firma muss für
die Mitarbeiter eine Gruppeneinteilung in eine gegebene Anzahl Gruppen
vornehmen.
Da sich das Management Sorgen darüber macht, dass unnötiges Palaver
die Kosteneffizienz des Events gefährden könnte, soll die Gruppeneinteilung
so gestaltet sein, dass die Gruppenmitglieder paarweise bereits in mindestens
einem Projekt zusammengearbeitet haben, sich also gegenseitig kennen.
Damit können die endlosen Vorstellungsrunden zu Beginn des Events
eingespart werden und es kann gleich mit der produktiven Arbeit begonnen
werden, so die Argumentation der Erbsenzähler im Management.
Nach längerem Nachdenken fragt sich der Organisator frustriert, warum
es so schwierig ist, eine solche Gruppeneinteilung zu finden.

\begin{loesung}
Das Problem \textit{GRUPPENEINTEILUNG} ist äquivalent zum Problem
\textit{CLIQUE-COVER}, wie die folgende Eins-zu-Eins-Reduktion
zeigt.
Die Mitarbeiter der Firma bilden die Knoten eines Graphen, eine
Kante bedeutet, dass die beiden Mitarbeiter bereits einmal zusammengearbeitet
haben.
Eine Gruppe soll aus Mitarbeitern bestehen, die sich alle kennen,
sie bilden eine Clique.
Die Gruppeneinteilung verlangt, dass $k$ Cliques gefunden werden 
müssen, so dass jeder Mitarbeiter in genau einer Clique ist.
\begin{center}
\begin{tabular}{r>{$}c<{$}l}
\textit{GRUPPENEINTEILUNG}&\leftrightarrow&\textit{CLIQUE-COVER}\\
\hline
Mitarbeiter&\leftrightarrow&Knoten\\
frühere Zusammenarbeit&\leftrightarrow&Kanten\\
Gruppe&\leftrightarrow&Clique\\
Anzahl $k$ der Gruppen&\leftrightarrow&Anzahl $k$ der Cliques\\
Jeder Mitarbeiter in genau einer Gruppe&\leftrightarrow&Jeder Knoten in genau einer CLIQUE
\end{tabular}
\end{center}
Da \textit{CLIQUE-COVER} NP-vollständig ist, ist auch
\textit{GRUPPENEINTEILUNG} NP-vollständig und es gibt nach aktuellem
Wissen keinen polynomiellen Algorithmus zur Bestimmung einer Gruppeneinteilung
nach den spezifizierten Regeln.
Jeder bekannte Algorithmus skaliert exponentiell mit der Grösse der
Firma, was zu den beobachteten langen Zeiten für die Lösung führt.
\end{loesung}

\begin{bewertung}
Reduktionsprinzip ({\bf R}) 1 Punkt,
Wahl eines geeigneten Vergleichsprobleme wählen ({\bf V}) 1 Punkt,
Abbildung für Mitarbeiter ({\bf M}) 1 Punkt,
Zusammenarbeit ({\bf Z}) 1 Punkt, und Gruppen ({\bf G}) 1 Punkt,
NP-Vollständigkeit ({\bf N}) 1 Punkt.
\end{bewertung}

